\documentclass{article}
% A leanblueprint-style source: each item carries a \label (its identity),
% \lean{...} linking it to its Lean declaration(s), \uses{...} dependencies,
% and \leanok once the Lean side checks. `hgraph sync` parses exactly these
% macros to build the graph — the blueprint is the reference, not a foil.
\usepackage{leanblueprint}
\begin{document}

\section{Gauss summation}

\begin{definition}[Even]\label{def:even}\lean{Gauss.IsEven}\leanok
A natural number $n$ is \emph{even} if $n = 2k$ for some $k$.
\end{definition}

\begin{lemma}[Sum of evens]\label{lem:even-add}\lean{Gauss.isEven_add}\uses{def:even}\leanok
If $a$ and $b$ are even then $a + b$ is even.
\end{lemma}

\begin{lemma}[Product with successor]\label{lem:mul-succ}\lean{Gauss.isEven_mul_succ}\uses{def:even}
For every $n$, the product $n(n+1)$ is even.
\end{lemma}

% This one is already in Mathlib — \mathlibok + a \lean pointing straight at the
% library declaration, so it needs no local formalization.
\begin{lemma}[Sum splits off the last term]\label{lem:sum-succ}\lean{Finset.sum_range_succ}\mathlibok
For every $n$, $\sum_{i=0}^{n+1} i = \bigl(\sum_{i=0}^{n} i\bigr) + (n+1)$.
\end{lemma}

\begin{theorem}[Gauss]\label{thm:gauss}\lean{Gauss.sum_range_id}
$2\sum_{i=0}^{n} i = n(n+1)$.
\end{theorem}
\begin{proof}\uses{lem:sum-succ, lem:mul-succ, lem:even-add}
By induction on $n$, splitting off the last term with the sum-splitting lemma;
evenness of $n(n+1)$ records the divisibility along the way.
\end{proof}

\end{document}
